\chapter*{Notación}
\addcontentsline{toc}{chapter}{Notación}

\begingroup
\small
\setlength{\LTleft}{0pt}
\setlength{\LTright}{0pt}
\renewcommand{\arraystretch}{1.15}

\section*{Notación matemática y algebraica}
\begin{longtable}{p{0.33\textwidth}p{0.61\textwidth}}
\toprule
\textbf{Símbolo} & \textbf{Significado} \\
\midrule
\endfirsthead
\toprule
\textbf{Símbolo} & \textbf{Significado} \\
\midrule
\endhead
$\F, \K$ & Cuerpo (finito, base, escalar), según el contexto. \\
$\Fp, \Fq, \Fr$ & Cuerpos finitos de orden $p$, $q$ y $r$. \\
$\charac{\K}$ & Característica del cuerpo $\K$. \\
$\GG, \GOne, \GTwo, \GT$ & Grupos algebraicos usados en esquemas con pairings. \\
$\pairing{P}{Q}$ & Pairing bilineal evaluado en $P \in \GOne$ y $Q \in \GTwo$. \\
$\E, \OO$ & Curva elíptica y su punto al infinito. \\
$\Hash$ & Función de hash genérica. \\
$\Prob$ & Operador de probabilidad. \\
$\securityparam, \negl$ & Parámetro de seguridad y función negligible. \\
$\BigO, \BigOmega, \BigTheta$ & Notación asintótica estándar. \\
$\Pclass, \NP, \IP, \PSPACE$ & Clases de complejidad computacional. \\
\bottomrule
\end{longtable}

\section*{Pruebas sucintas y aritmetización}
\begin{longtable}{p{0.33\textwidth}p{0.61\textwidth}}
\toprule
\textbf{Símbolo} & \textbf{Significado} \\
\midrule
\endfirsthead
\toprule
\textbf{Símbolo} & \textbf{Significado} \\
\midrule
\endhead
$\Prover, \Verifier$ & \textit{Prover} y \textit{verifier} de un protocolo. \\
$\Relation, \Language$ & Relación y lenguaje asociados a un problema de verificación. \\
$\InstanceSpace, \WitnessSpace, \ProofSpace$ & Espacios de \textit{public inputs}, \textit{witness} y pruebas/argumentos. \\
\Setup, \Prove, \Verify & Algoritmos de protocolos de pruebas/argumentos. \\
$\CRS$, \SRS & \textit{Common Reference String} y \textit{Structured Reference String}. \\
\NIZK, \SNARK, \zkSNARK, \STARK, \zkSTARK & Familias de pruebas no interactivas y sucintas. \\
\ROneCS, \QAP, \PCS, \VCS & Modelos y primitivas de aritmetización y compromiso. \\
\KZG, \FiatShamir, \Groth, \PLONK, \HaloTwo & Protocolos usados recurrentemente en la tesis. \\
$\Circuit, \Constraint$ & Circuito aritmético y restricción genérica de un circuito. \\
$\EvalDomain, \VanishingPoly$ & Dominio de evaluación y polinomio anulador asociado. \\
$\QuotientPoly, \PermutationPoly, \LookupPoly$ & Polinomio cociente y polinomios auxiliares de permutación y \textit{lookup}. \\
\bottomrule
\end{longtable}

\section*{Merkle Trees}
\begin{longtable}{p{0.33\textwidth}p{0.61\textwidth}}
\toprule
\textbf{Símbolo} & \textbf{Significado} \\
\midrule
\endfirsthead
\toprule
\textbf{Símbolo} & \textbf{Significado} \\
\midrule
\endhead
$\MerkleTree{T}$ & Árbol de Merkle genérico. \\
$\MerkleRoot{\MerkleTree{T}}$ & Raíz del árbol $\MerkleTree{T}$. \\
$\MerkleProof$ & Prueba de inclusión de Merkle. \\
$\MerklePath$ & Camino autenticado dentro del árbol. \\
$\MerkleLeaf, \MerkleSibling, \MerkleIndex$ & Hoja, nodo hermano e índice de una prueba de inclusión. \\
\SMT & \textit{Sparse Merkle Tree}. \\
$\EmptyNode$ & Hoja o nodo vacío en un \SMT. \\
$\RootIn, \RootOut$ & Raíces de entrada y salida del anti-replay. \\
$\SiblingVec$ & Vector de nodos hermanos en una prueba de Merkle. \\
\bottomrule
\end{longtable}

\section*{\Cardano y validadores}
\begin{longtable}{p{0.33\textwidth}p{0.61\textwidth}}
\toprule
\textbf{Símbolo} & \textbf{Significado} \\
\midrule
\endfirsthead
\toprule
\textbf{Símbolo} & \textbf{Significado} \\
\midrule
\endhead
$\ChainA, \ChainB$ & Cadenas de origen y destino del bridge. \\
$\state, \tx, \event$ & Estado, transacción y evento abstractos del protocolo. \\
\LockTx, \MintTx, \BurnTx, \UnlockTx, \UpdaterTx & Transacciones canónicas del flujo del bridge. \\
\UTxO, \UTxOs, \eUTxO, \eUTxOs & Modelo \UTxO clásico y su extensión en \Cardano. \\
\Plutus, \UPLC & Lenguaje y \textit{assembly} de scripts en \Cardano. \\
$\Ouroboros$ & Protocolo de consenso de \Cardano. \\
$\Epoch, \Slot$ & Época y slot del protocolo \Ouroboros. \\
$\datum, \redeemer, \validator, \txcontext$ & Componentes semánticos de un validador en \Cardano. \\
\bottomrule
\end{longtable}

\section*{\Mithril, certificados y snapshots}
\begin{longtable}{p{0.33\textwidth}p{0.61\textwidth}}
\toprule
\textbf{Símbolo} & \textbf{Significado} \\
\midrule
\endfirsthead
\toprule
\textbf{Símbolo} & \textbf{Significado} \\
\midrule
\endhead
\Mithril, \STM & Protocolo \Mithril y su esquema de multi-firma. \\
$\Stake, \StakeDistribution$ & Participación individual y distribución de stake certificada. \\
\AVK & \textit{Aggregate Verification Key} del certificado vigente. \\
$\ProtocolMsg, \Certificate$ & Mensaje de protocolo certificado y certificado genérico. \\
\GenesisCertificate, \StakeDistributionCertificate, \TransactionSnapshotCertificate & Certificados de \Mithril: génesis de la cadena de certificados, de \StakeDistribution y de \textit{snapshot} de transacciones. \\
\TransactionSnapshot & Snapshot de transacciones expuesto por \Mithril. \\
$\TransactionSnapshotHash$ & Hash asociado a un snapshot de transacciones. \\
\CertificateChain & Cadena de certificados enlazados. \\
$\ChainVerify, \CertVerify$ & Verificación de la cadena completa de \Mithril. \\
$\CertVerify$ & Verificación de un certificado individual de \Mithril. \\
\STMRelation, \STMLanguage & Relación y lenguaje del circuito \STM. \\
\TxSetUpdateRelation & Relación del circuito del conjunto de transacciones usadas. \\
\TxSnapshotMembershipRelation & Relación del circuito que prueba pertenencia a un snapshot. \\
\bottomrule
\end{longtable}

\endgroup
